%% =========================================================================
%%  Обозначения и сокращения
%%  (наименование структурного элемента по Положению о ВКР МФТИ, Прил. 1 §3)
%% =========================================================================

\chapter*{Обозначения и сокращения}
\addcontentsline{toc}{chapter}{Обозначения и сокращения}

\section*{Базовые объекты}
\addcontentsline{toc}{section}{Базовые объекты}

\begin{description}[leftmargin=5.5em,labelwidth=5em,labelindent=0pt,font=\normalfont]
  \item[$F$]       отображение $F\colon \RR^n \to \RR^n$ в задаче $F(x)=0$
                   (главы~1,~2) или градиент $F = \nabla f$ задачи оптимизации (глава~3)
  \item[$f$]       скалярная целевая функция $f\colon\RR^n\to\RR$ в оптимизации
  \item[$J(x)$]    якобиан $F$ в точке $x$; в задаче оптимизации~--- гессиан $\nabla^2 f(x)$
  \item[$J^*$]     якобиан в решении: $J^* = J(x^*) = \nabla F(x^*)$
  \item[$x^*$]     решение: $F(x^*)=0$
\end{description}

\section*{Квазиньютоновские обновления}
\addcontentsline{toc}{section}{Квазиньютоновские обновления}

\begin{description}[leftmargin=5.5em,labelwidth=5em,labelindent=0pt,font=\normalfont]
  \item[$x_k$]      приближение к решению на $k$-й итерации
  \item[$s_k$]      шаг: $s_k = x_{k+1} - x_k$
  \item[$y_k$]      разность значений: $y_k = F(x_{k+1}) - F(x_k)$
  \item[$d_k$]      направление поиска: $B_k\,d_k = -F(x_k)$
  \item[$B_k$]      текущая аппроксимация якобиана
  \item[$\widehat{B}_k$] базовая матрица в формуле обновления
                    (накопление: $\widehat{B}_k = B_k$; пересборка:
                    $\widehat{B}_k = \widetilde{B}$, обычно $I$)
  \item[$\widetilde{B}$] опорная матрица~--- параметр алгоритма
                    SP-Broyden, не зависящий от $k$;
                    к ней производится сброс при пересборке
  \item[$\Delk$]    изменение матрицы: $\Delk = B_{k+1} - \widehat{B}_k$
  \item[$S_m,\, Y_m,\, V_m$]
                    матрицы прошлых шагов, разностей и направлений
                    обновления в общей мультисекущей формуле
                    (\ref{eq:general}); $V_m$ ($m$~--- ранг обновления)
                    параметризует семейство
  \item[$E_k$]      ошибка аппроксимации якобиана: $E_k = B_k - J^*$
\end{description}

\section*{SP-Broyden и SS-U}
\addcontentsline{toc}{section}{SP-Broyden и SS-U}

\begin{description}[leftmargin=5.5em,labelwidth=5em,labelindent=0pt,font=\normalfont]
  \item[$m$]           ранг мультисекущего обновления в общей формуле
                        главы~\ref{ch:jacobian} ($m$ секущих условий
                        $B_{k+1} S_m = Y_m$); используется только
                        в главе~\ref{ch:jacobian}
  \item[$p$]           глубина сохраняемых \emph{прошлых} секущих условий
                        в SP-Broyden и его вариантах (главы~\ref{ch:sp-broyden-theory},
                        \ref{ch:symmetric}); число столбцов окна на одну
                        больше, чем $p$ --- $S_p \in \RR^{n\times (p+1)}$
                        за счёт включения текущей секущей $s_k$
  \item[$S_p,\, Y_p$]  окно текущего и $p$ прошлых шагов и разностей,
                        $S_p = [s_k\mid s_{k-1}\mid\cdots\mid s_{k-p}]
                        \in \RR^{n\times(p+1)}$
  \item[$G_p$]         грамиан окна: $G_p = S_p^\top S_p$
  \item[$v_k^*$]       оптимальный SP-вектор:
                        $v_k^* = S_p\,G_p^{-1}\,e_1$
  \item[$\Pi_k$]       ортогональный проектор на $\col(S_p)$:
                        $\Pi_k = S_p\,G_p^{-1}\,S_p^\top$;
                        $\Pi_k^{\perp} = I - \Pi_k$
  \item[$\tau$]        порог числа обусловленности грамиана
                        (адаптивный выбор глубины $p$)
  \item[$u_k^*$]       ведущий собственный вектор коррекции SS-U
                        (в $s_k^{\perp}$)
  \item[$\sigma_k^*$]  коэффициент ранг-1 коррекции SS-U
  \item[$\mathcal{U}$] базовое симметричное квазиньютоновское обновление
                        (SR1, PSB) в конструкции SS-U
\end{description}

\section*{Соглашения}
\addcontentsline{toc}{section}{Соглашения}

\begin{itemize}[leftmargin=*]
  \item Индекс $k$ нумерует итерации алгоритма;
        индексы внутри окна пробегают $j = k-p,\ldots,k$.
  \item $I$~--- единичная матрица подходящей размерности;
        $e_1 = (1,0,\ldots,0)^\top$~--- первый базисный вектор.
  \item Если контекст не указан явно, все нормы матриц~--- фробениусовы
        ($\Fro{\cdot}$), все нормы векторов~--- евклидовы ($\|\cdot\|$).
  \item Константы обозначаются прописными латинскими буквами
        ($L,\, C_d$); параметры алгоритма~---
        строчными греческими ($\tau,\, \bar{\rho}$);
        индексы $k$ указывают итерационную зависимость.
\end{itemize}
